% from ega4-11-fr.tex
\oldpage[IV-3]{161}

On dit que la famille $(u_\lambda)$ (ou la famille correspondante des
$u^\#_\lambda:(f_\lambda)^*(\mathcal F)\to\mathcal G_\lambda$) est
\emph{séparante} si l'intersection des noyaux des $u_\lambda$ est
\emph{nulle}. En d'autres termes, cela signifie que pour \emph{tout} ouvert
$U$ de $X$, et toute section $t$ de $\mathcal F$ au-dessus de $U$, telle que,
pour \emph{tout} $\lambda$, la section $u_\lambda(t)$ (qui, par définition est
une section de $\mathcal G_\lambda$ au-dessus de $f_\lambda^{-1}(U)$) soit
nulle, alors $t$ est elle-même nulle.

\begin{env}[11.9.2]
\label{IV.11.9.2-fr}
Avec les notations de (11.9.1), soit $M$ un second ensemble d'indices~; pour
tout $\lambda\in L$, soit $(g_{\lambda\mu})_{\mu\in M}$ une famille de
morphismes $g_{\lambda\mu}:Z'_{\lambda\mu}\to Z_\lambda$~; pour tout couple
$(\lambda,\mu)$, supposons donnés un $\mathcal O_{Z'_{\lambda\mu}}$-Module
quasi-cohérent $\mathcal H_{\lambda\mu}$ et un homomorphisme
$v_{\lambda\mu}:\mathcal G_\lambda\to
(g_{\lambda\mu})_*(\mathcal H_{\lambda\mu})$~; posons
$h_{\lambda\mu}=f_\lambda\circ g_{\lambda\mu}:Z'_{\lambda\mu}\to X$ et
considérons l'homomorphisme composé
\[
w_{\lambda\mu}:\mathcal F\xrightarrow{u_\lambda}
(f_\lambda)_*(\mathcal G_\lambda)
\xrightarrow{(f_\lambda)_*(v_{\lambda\mu})}
(h_{\lambda\mu})_*(\mathcal H_{\lambda\mu}).
\]
Supposons que, pour tout $\lambda\in L$, la famille
$(v_{\lambda\mu})_{\mu\in M}$ soit séparante~: alors, il en est de même de la
famille des $(f_\lambda)_*(v_{\lambda\mu})$ ($\mu\in M$), comme on le voit
aussitôt. On en conclut que, pour que la famille $(u_\lambda)$ soit séparante,
il faut et il suffit que la famille
$(w_{\lambda\mu})_{(\lambda,\mu)\in L\times M}$ le soit.
\end{env}

\begin{env}[11.9.3]
\label{IV.11.9.3-fr}
Pour vérifier que la famille $(u_\lambda)$ est séparante (avec les notations de
(11.9.1)), on peut évidemment se ramener tout d'abord au cas où $X$ est
\emph{affine}, la propriété étant locale sur $X$. On peut en outre supposer que
$Z_\lambda=X$ \emph{pour tout} $\lambda\in L$. En effet, soit $M$ un ensemble
d'indices, somme d'une famille $(M_\lambda)_{\lambda\in L}$, et pour tout
$\lambda\in L$, soit $(Y_{\lambda\mu})_{\mu\in M_\lambda}$ un recouvrement
ouvert affine de $Z_\lambda$~; soit
$j_{\lambda\mu}:Y_{\lambda\mu}\to Z_\lambda$ l'injection canonique et posons
$\mathcal H_{\lambda\mu}=j_{\lambda\mu}^*(\mathcal G_\lambda)
=\mathcal G_\lambda|Y_{\lambda\mu}$. Si l'on considère l'homomorphisme
canonique
$v_{\lambda\mu}=\rho_{\mathcal G_\lambda}:\mathcal G_\lambda\to
(j_{\lambda\mu})_*(j_{\lambda\mu}^*(\mathcal G_\lambda))
=(j_{\lambda\mu})_*(\mathcal H_{\lambda\mu})$ relatif à $j_{\lambda\mu}$
$(\mathbf 0_{\mathrm I},4.4.3.2)$, il est immédiat que pour chaque
$\lambda\in L$, la famille $(v_{\lambda\mu})_{\mu\in M_\lambda}$ est
séparante. En vertu de (11.9.2), on est donc ramené à prouver que la famille des
homomorphismes composés $((f_\lambda)_*(v_{\lambda\mu}))\circ u_\lambda$ est
séparante, autrement dit on est ramené au cas où les $Z_\lambda$ sont affines.
Mais alors les $(f_\lambda)_*(\mathcal G_\lambda)$ sont des
$\mathcal O_X$-Modules quasi-cohérents (\textbf I, 1.6.2) et en vertu de la
définition, on peut remplacer les $Z_\lambda$ par $X$ et les
$\mathcal G_\lambda$ par les $(f_\lambda)_*(\mathcal G_\lambda)$, d'où notre
assertion.

On notera en outre que si $L$ est \emph{fini} et les $f_\lambda$
\emph{quasi-compacts}, on peut, dans la réduction précédente, supposer que les
$M_\lambda$ sont aussi \emph{finis}, donc on est dans ce cas ramené à vérifier
qu'une famille \emph{finie} d'homomorphismes de $\mathcal F$ dans des
$\mathcal O_X$-Modules quasi-cohérents est séparante.
\end{env}

\begin{env}[11.9.4]
\label{IV.11.9.4-fr}
Considérons donc le cas où $Z_\lambda=X$ pour tout $\lambda$, et où
$X=\operatorname{Spec}(A)$ est affine~; alors on a
$\mathcal F=\widetilde M$ et $\mathcal G_\lambda=\widetilde N_\lambda$, où
$M$ et $N_\lambda$ sont des $A$-modules, et $u_\lambda=\widetilde\varphi_\lambda$,
où les $\varphi_\lambda:M\to N_\lambda$ sont des $A$-homomorphismes. Dire que
la famille $(u_\lambda)$ est séparante signifie alors que, pour tout $s\in A$,
l'intersection des noyaux des
$(\varphi_\lambda)_s:M_s\to(N_\lambda)_s$ est réduite à 0. On dit alors aussi
que la famille $(\varphi_\lambda)$ est \emph{séparante}. On notera que si $L$
est \emph{fini}, il revient au même de dire que l'intersection des noyaux des
$\varphi_\lambda$ est 0, car on a alors
\[
\bigcap_{\lambda\in L}\operatorname{Ker}((\varphi_\lambda)_s)
=\left(\bigcap_{\lambda\in L}\operatorname{Ker}(\varphi_\lambda)\right)_s
\quad(\mathbf 0_{\mathrm I},1.3.2).
\]
Mais cette relation n'est plus exacte en général lorsque $L$ est \emph{infini},
et le fait que l'intersection des noyaux des $\varphi_\lambda$ soit 0
\emph{n'entraîne donc pas},
